%==============================================================================
\section{Mean, Median, Mode --- Trung bình, trung vị, mode}
\label{sec:mean-median-mode}
%==============================================================================

\begin{dinhnghia}[title={Ba đo xu hướng trung tâm}]
Mean (trung bình), Median (trung vị) và Mode (mode) là các đại lượng mô tả \textbf{xu hướng trung tâm} của dataset.
Trong ML, chúng giúp hiểu phân phối dữ liệu và nhận diện outlier.
\end{dinhnghia}

\subsection{Mean (Trung bình)}

\begin{ynghia}[title={Mean}]
Mean là \textbf{giá trị trung bình}: cộng tất cả giá trị rồi chia cho số quan sát.
Mean nhạy với \textbf{outlier} --- giá trị cực đoan kéo mean mạnh.
\end{ynghia}

\begin{phuongphap}[title={Python}]
Dùng \texttt{np.mean()} trong NumPy.
\end{phuongphap}

\subsection{Median (Trung vị)}

\begin{ynghia}[title={Median}]
Median là giá trị \textbf{ở giữa} sau khi sắp xếp.
Nếu số phần tử chẵn, median là trung bình của hai giá trị giữa.
Median \textbf{ít bị ảnh hưởng} bởi outlier hơn mean.
\end{ynghia}

\begin{phuongphap}[title={Python}]
Dùng \texttt{np.median()}.
\end{phuongphap}

\subsection{Mode (Mode)}

\begin{ynghia}[title={Mode}]
Mode là giá trị \textbf{xuất hiện nhiều nhất}.
Nếu nhiều giá trị cùng tần suất cao nhất, dataset có thể lưỡng cực (bimodal), tam cực (trimodal) hoặc đa cực (multimodal).
Mode hữu ích khi tìm giá trị phổ biến nhất; kém phù hợp khi dữ liệu trải rộng hoặc không có giá trị lặp.
\end{ynghia}

\begin{phuongphap}[title={Python}]
SciPy có \texttt{mode()}; Pandas: \texttt{Series.mode()}.
\end{phuongphap}

\subsection{Triển khai Python}

\begin{vidu}[title={Bảng lương mẫu}]
\begin{lstlisting}
import numpy as np
import pandas as pd
# create a sample salary table
salary = pd.DataFrame({
   'employee_id': ['001', '002', '003', '004', '005', '006', '007',
   '008', '009', '010'],
   'salary': [50000, 65000, 55000, 45000, 70000, 60000, 55000, 45000,
   80000, 70000]
})

# calculate mean
mean_salary = np.mean(salary['salary'])
print('Mean salary:', mean_salary)

# calculate median
median_salary = np.median(salary['salary'])
print('Median salary:', median_salary)

# calculate mode
mode_salary = salary['salary'].mode()[0]
print('Mode salary:', mode_salary)
\end{lstlisting}

\textbf{Output}:
\begin{lstlisting}
Mean salary: 59500.0
Median salary: 57500.0
Mode salary: 45000
\end{lstlisting}
\end{vidu}

\begin{ynghia}[title={Đọc kết quả}]
Mean cao hơn median một chút (ảnh hưởng lương 70k--80k); mode 45000 vì hai nhân viên cùng mức lương đó.
\end{ynghia}
